\chapter{方向不对称：正向生成、逆向辨识与路径预算}
\label{chap:path_integration_asymmetry}

上一章把方向性任务代价与更新顺序效应分开定义，本章进一步研究前一问题：给定相同的两个表征端点，为什么从 $A$ 到 $B$ 与从 $B$ 到 $A$ 往往不是同一计算过程的简单倒放。关键不在“思维逆着自然熵流做功”，而在于两个任务拥有不同的输入信息、候选集合、损失函数、可辨识条件和验收程序。我们将对称的表示距离与方向相关的执行成本分账，用路径泛函记录搜索、验证、通信和风险，再分析正向生成、逆向推断以及两者之间的回路差异。原有顺流、逆流和磁滞图像仍可帮助定位现象，但结论必须落到算法、证据、资源、离散实现及外部验收上，不能把统计代价、信息熵、活动强度和焦耳能耗视作同一量。

\section{双路径模型：表示距离与方向任务成本}
\label{sec:dual_path_model}

设任务相关状态空间为 $\mathcal Z_\tau$，其中同时包含全局分布表示 $h\in\mathbb R^{d_h}$、带身份和版本的局部结构网络 $m$，以及可用证据 $e$。对两个状态 $z_A,z_B$，我们可以定义无量纲的对称表示距离
\[
d_G(A,B)=\inf_{\gamma:A\to B}\int_0^1
\sqrt{\dot\gamma(s)^\top G(\gamma(s),\tau)\dot\gamma(s)}\,ds,
\]
其中 $G$ 是正定度量，曲线 $\gamma$ 在固定结构版本和任务条件下分段可微。由此有 $d_G(A,B)=d_G(B,A)$，但它只说明在选择的特征与度量下两个状态相隔多远，不说明从一个状态计算另一个状态需要多少时间，也不保证相近状态具有相同身份、因果关系或行动后果。若结构网络发生离散编辑，距离还要补充节点身份、边类型和版本的匹配代价；若度量、目标或输入随时间变化，前后两次距离也不能直接比较。

方向任务成本则以过程为对象。令 $\gamma=(z_0,a_0,u_0,\ldots,z_N)$ 是从输入状态到验收集合的离散轨迹，步长为 $\Delta t_k>0$，运行成本率、验证成本、通信成本和风险分别记作 $c_r,c_v,c_q,\rho$。给定任务 $\tau$ 和载体 $\chi$，定义
\[
\mathcal J_\tau(\gamma)=\sum_{k=0}^{N-1}\Delta t_k
[c_r(z_k,a_k)+c_v(z_k)+c_q(z_k)+\lambda_\rho\rho(z_k,a_k)]
+\Phi_\tau(z_N),
\]
其中终端项 $\Phi_\tau$ 衡量未达到验收标准的残差。各项只有在单位归一化且权重由任务方声明时才能相加，否则应保留资源向量和硬约束。$A\to B$ 与 $B\to A$ 的输入可见信息、动作集合与终端集合通常不同，因此最优轨迹 $\gamma^*_{A\to B}$、$\gamma^*_{B\to A}$ 及其成本可以不对称。这种不对称来自任务定义和可用算法，而不是对称距离自身偷偷变成“有方向的路程”。

路径积分语言仍有用，但必须说明其资格。连续模型可把上式写成 $\int L_\tau(z,\dot z,a,u,t)dt$；只有 $L$ 的对象、量纲和边界给定后，极值路径才有计算含义。若工具调用、数据库写入或权限提交造成跳变，轨迹应由连续段和事件重置 $z^+=R_e(z^-)$ 共同构成，不能用一条光滑曲线掩盖不可逆事务。所谓“事实距离”也不是本体论距离或共现概率的天然同义词，它只是任务相关表征空间中的模型量。我们通过保留查询、身份核验和扰动实验检查它是否适合当前任务。

双路径图据此改写如下。绿线表示给定生成接口的低成本候选轨迹，红线表示给定反演接口的搜索—验证轨迹；两线围成面积只是历史状态、可用证据和操作集合不同的可视化，不证明存在认知磁场或热力学磁滞。若生成过程保存了充分日志，逆向轨迹可能缩短；若生成结果要满足严格形式约束，绿线也可能比红线更长。

\begin{figure}[htbp]
    \centering
    \begin{adjustbox}{max width=0.98\linewidth,max totalheight=0.72\textheight}
\begin{tikzpicture}
        % 定义点 A 和 B
        \node (A) at (0,0) [circle, fill=blue!20, draw=blue, thick] {A};
        \node (B) at (6,0) [circle, fill=red!20, draw=red, thick] {B};

        % 顺流路径 (弧线向下)
        \draw[->, thick, green!70!black] (A) to [bend right=45] node[midway, below] {$\gamma_{gen}$ (顺流/生成/低 Cost)} (B);

        % 逆流路径 (弧线向上，更长，更曲折)
        \draw[->, thick, red!70!black] (B) to [out=135, in=45, looseness=1.5] node[midway, above] {$\gamma_{inf}$ (逆流/推理/高 Cost)} (A);

        % 规范场风向
        \draw[->, dashed, gray] (1, -1) -- (5, -1) node[right] {$\vec{E}$ (规范场/热力学箭头)};
    \end{tikzpicture}
\end{adjustbox}
    \caption{双路径模型：顺流的生成与逆流的推理形成了一个非零面积的回路（磁滞环）。}
\end{figure}

因此，本节的模型假设是状态与任务可记录、局部度量可估计、资源账本可比较；数学结果是给定这些对象后可以分别优化两条轨迹；经验主张则是某个系统在某类数据上确有方向差异，必须由配对实验确认。三者不能互相替代。局部结构网络负责端点身份、来源和版本，全局分布表示负责连续候选及置信度，两者由任务绑定和读出接口联合，而非由图形相似推出同构。

配对实验还要控制训练暴露与缓存命中，否则方向差异可能只是数据泄漏。我们对同一端点对同时构造生成任务和辨识任务，固定成功阈值、载体、可用工具与截止时间，分别报告条件成功率下的中位数和尾部成本；随后交换提示顺序、隐藏日志并扰动候选规模。如果不对称只在熟悉样本上出现，它更可能来自记忆索引；如果在候选数增长时呈稳定复杂度差异，才支持算法解释。即便如此，结论也只适用于声明的问题族，不能把局部测量升格为所有认知过程的普遍定律。

\section{正向路径：条件生成、默认转移与约束验收}
\label{sec:forward_flow}

正向路径从给定原因、条件或前缀产生结果。设输入为 $x$，生成状态为 $h_k$，输出符号为 $y_k$，模型按
\[
h_{k+1}=F_\theta(h_k,x,y_k,u_k),\qquad
y_{k+1}\sim p_\theta(\cdot\mid h_{k+1},x)
\]
递推。$F_\theta$ 的参数和结构在本次运行中固定，采样温度、搜索宽度、停止条件和工具反馈进入控制量。这个过程之所以常显得“顺流”，是因为训练已把常见转移压缩到参数和索引中，运行时只需沿已学习条件分布展开；此前付出的数据整理、训练、对齐和缓存成本被摊销到了载体，而不是消失。一次生成的低延迟不能证明全过程低成本，更不能证明它符合热力学熵增或设备能耗极低。

“看到乌云立即想到下雨”是默认转移的例子。系统可由历史共现迅速提高“下雨”候选权重，但乌云也可能消散，地面天气还受季节、风场和地点影响。直觉读出只是一条高先验候选，不是因果证明。工程上应把来源、地点和时间绑定到局部结构记录，把降雨概率和备选解释保存在分布状态中；行动前再依据风险决定是否调用天气工具。若任务只是补全一句话，快速候选可能足够；若任务是发出洪水警报，则需要独立观测、校准和权限检查。相同内部联想因验收标准不同而具有不同任务成本。

LLM 的下一个 token 生成也不是“顺概率梯度滑行”的完整描述。模型先计算条件 logits，再由解码规则选择或采样；贪心、束搜索、核采样与受约束解码的路径和成本均不同。生成一个语法通顺的句子可以很快，生成一段满足接口 schema、引用真实、可执行且不过权的计划则需要验证、检索和重试。概率高只表示模型分布中的相对偏好，不等于事实正确、价值正当或环境可行。我们因此把生成成本拆为候选展开、约束过滤、证据核验和提交回执，而不把所有正向过程归为低难度。

信息扩散同样要限定。由一个精确输入产生多个可能输出，会增加候选分布的熵；但确定性程序也能从原因计算唯一结果，受约束生成还可能使输出分布更集中。Shannon 熵描述指定随机变量的分布，不能直接充当运行代价、理解程度或焦耳能耗。若要比较前后不确定性，必须固定样本空间和条件；候选集合变化时要重新定义分布。设备能耗则由功率计或载体模型估计，不能从输出熵的正负推出。

正向路径的验收至少包括四项：第一，输出是否达到任务损失与格式约束；第二，身份、来源与版本是否一致；第三，事实与环境后果是否经独立接口核验；第四，运行资源与风险是否在预算内。AgentOS 示例中，SPC 可以压缩当前候选，TCE 调整搜索强度，CCM 监测 $h(t)$ 占用，Boundary 限制外部动作，Offloading 把验证交给合适工具，$M_e$ 保存可重放证据。这些模块只是一种实现，不应被规定为所有生成系统的唯一机制。

反例可以检验“正向必然容易”的说法。生成满足数百项约束的工程设计、形式证明或多主体承诺，可能比验证一个给定候选更难；分布外输入会让默认转移失效；工具延迟、上下文拥塞与身份冲突也会把短路径变成长路径。实验应比较相同成功率下的分位时延、计算量、工具调用和修复次数，而不能只选择成功样本的平均响应时间。正向路径的优势是可被训练、缓存和接口复用，不是宇宙赋予的无条件下坡。

\section{逆向路径：可辨识性、搜索验证与学习更新}
\label{sec:reverse_flow}

逆向任务从观测结果 $y$ 推断潜在原因 $x$。先定义生成模型 $p(y\mid x,m)$、先验 $p(x\mid m)$ 和结构版本 $m$，则后验满足
\[
p(x\mid y,m)=\frac{p(y\mid x,m)p(x\mid m)}{\int p(y\mid x',m)p(x'\mid m)dx'}.
\]
这一定义显示逆向困难首先来自多对一映射、噪声和先验不充分，而不是“逆熵攀登”。若不同原因给出近似相同观测，问题不可辨识；再多计算也不能从现有证据恢复唯一真相。系统应输出候选集合、后验权重和还需采集的辨别证据，而不是把搜索到的单一原因包装成确定结论。

“地湿了，是否刚下雨”正好说明这一点。下雨、洒水车、管道泄漏和人工清洗都能产生湿地。逆向过程可先建立候选因果图，再查询云量、降雨记录、水痕分布和车辆日志。每次查询会改变后验，也消耗时延、通信和权限预算。若系统只沿最高先验选择“下雨”，它完成的是快速猜测而非可靠推断；若通过干预或多源证据排除其他解释，才提高辨识度。知识结构的价值在于组织候选和证据依赖，不在于为推理注入神秘“第三驱动力”。

单向函数与哈希展示更强边界。已知 $y=f(x)$ 时，验证候选 $x'$ 是否满足 $f(x')=y$ 可能容易，但从 $y$ 构造原像可能计算困难；若 $f$ 多对一，即使无限算力也不能确定原始输入而只能给原像集合。困难程度依安全参数、算法族、侧信道和硬件而变，不能从“小整数分解稍难”外推。所谓原路“断裂”，在计算上应明确为信息丢失、不可辨识或复杂度下界；若日志保存随机种子和中间状态，逆向可能显著变易。

逆向算法通常交替执行提出、评分与验证。设候选集为 $X_k$，提案核为 $q_k(x'\mid x,y)$，验证损失为 $\ell_v(x';y)$，则一步搜索可写作
\[
x_{k+1}\sim q_k(\cdot\mid x_k,y),\qquad
w_{k+1}\propto w_k\exp[-\beta_k\ell_v(x_{k+1};y)].
\]
$\beta_k$ 是无量纲选择强度，不是物理逆温度；步数、并行宽度、验证成本和停止置信度共同决定预算。蒙特卡洛、优化、符号求解和外部工具只是不同提案机制。所谓“升维绕行”若有实际含义，应表述为引入潜变量、中间证明、辅助观测或新的工具接口；不能用高维几何名称代替算法。

学习比一次逆推多出结构更新。新证据可能修改参数、身份边或规则版本，因而必须先在影子状态验证保留查询，再提交或回滚。连续参数更新声明步长与稳定条件，离散结构更新声明事务、权限和引用迁移。目标、输入、度量或结构变化时，总导数包含相应项，旧路径成本需要重新估计。AgentOS 可由 TCE 管理时间一致性、CCM 监控搜索复杂度、Boundary 控制证据与写入权限、Offloading 调用求解器，但它们不保证后验真实或学习成功。

逆向困难也有反例。带唯一条码的对象识别、附完整追踪信息的生成记录、线性可观测系统的状态估计，都可使恢复很直接；而正向生成复杂合规计划可能更难。验证必须在配对数据上报告成功率、校准、候选覆盖、查询数、时延和风险，并与只用先验、随机搜索及泄漏信息的负对照比较。只有在这些条件下，才能说某个逆向任务比对应正向任务更昂贵。

还需区分“找到一个可行解释”和“识别历史上的真实原因”。前者只要求候选通过观测约束，后者要求来源证据把候选绑定到实际事件；在欠定问题中，两者可能永远不能合并。例如多个程序生成同一输出时，求得任一程序是合成任务，确认当时执行的是哪一个程序则依赖日志、签名或见证。系统若缺少这些证据，应明确返回不可判定范围，并把主动采样的期望信息收益与采样风险一起交给策略层，而不是用更长的推理文本填补证据空缺。

\section{回路差异：不可逆记录、绕行策略与闭环验收}
\label{sec:section_4854}

把正向与逆向轨迹首尾相接会形成回路，但回路差异不能自动称为磁滞或做功。设正向算子为 $G$，逆向算子为 $I$，任务距离为 $d_\tau$，则重构残差定义为
\[
\varepsilon_{A}=d_\tau(I(G(A)),A),\qquad
\varepsilon_{B}=d_\tau(G(I(B)),B).
\]
残差来自压缩损失、随机采样、模型失配、证据缺失和结构写入。若 $G$ 可逆且 $I=G^{-1}$，两者可以为零；若提交改变外部环境或记忆版本，返回同一表征也不等于恢复同一历史状态。回路的可观测量应包括端点误差、状态差分、资源总量、失败次数和不可撤销事件，而不是未经定义的“理解功”。

绕行策略是逆问题中的正常设计。第一类是补充观测：为地湿案例查询天气、摄像和管网数据；第二类是改变表示：引入潜变量、因果图或中间证明，使候选更易区分；第三类是借助外部记忆和工具：索引、求解器、模拟器或其他 Agent 可以缩短本地搜索；第四类是降低要求：当唯一恢复不可行时，输出等价类、区间或风险上界。每种绕行都要把接口延迟、数据质量、权限和外部失败计入预算，不能把成本转移到环境后宣称系统内部实现了零代价理解。

我们把“理解”定义为任务相关能力集合，而不是回路积分。对给定问题族，系统若能压缩关键关系、生成可检验预测、解释反例、在扰动下迁移并根据证据修订，就获得较强理解证据。记住大量样本但无法辨别反例不足以证明理解；经历高成本搜索但输出错误也不构成理解。评价应包含保留集、分布外扰动、因果干预、校准和外部行动结果。内部损失下降、路径变长或资源消耗增加都不能单独替代成功。

离散闭环需要明确提交协议。每次正向生成保留输入、随机种子、模型版本和工具回执；逆向辨识保留候选、否证证据与停止原因；结构学习先创建候选版本 $m'$，通过身份一致性、权限、保留查询和回滚演练后再原子提交。若环境已被行动改变，则用补偿事务而非假装原路返回。连续模型的局部下降性质只有在步长、延迟和投影条件满足时才传给离散实现，结构切换还需单独检查驻留时间和重置稳定性。

全章验收分四组。对象验收检查状态、身份、任务、资源与能耗分账；算法验收比较 $G$、$I$、随机基线和带额外证据版本；动力学验收检查时变目标、输入、结构、度量和离散步长；任务验收由独立接口核对事实、约束、风险和环境结果。图中的两条路径只有在这些记录齐全时才是可复现实验对象，否则只是解释性插图。

因此，生成能力与反思能力确实是完整智能的重要组成，但二者不是“才华释放”和“功力注入”的物理二分。正向系统需要约束、证据与提交能力，逆向系统需要可辨识分析、搜索、验证与诚实的不确定性表达。HSF-HD在此提供一般状态动力学的共同接口，AgentOS提供一种工程落点；它们都不能从顺流、逆流、熵或做功的类比中自动获得正确性。只有路径对象可定义、资源账本可审计、反例可进入更新、外部结果可验收时，方向不对称才从诗意图像变成可用理论。

最终的关闭条件不是“两条曲线重新相交”，而是证据链、版本链与行动回执能够共同回答三个问题：系统产生了什么，依据什么把它反推回来，以及哪一部分损失已不可恢复。凡是不能回答的部分都应作为残余不确定性留在状态中，并影响后续预算和权限。这样，路径回路才成为可审计的计算闭环，而非用漂亮图形掩盖缺失记录。
